\clearpage
\VerdierPhysicalPageAnchor{0050}%
\VerdierSemanticAnchor{ch1-1-4-7-continuation}
其中 $(F_i,p_i)$ 分别由函子 $(F_i,m_i)$ 与截面 $s_i$ 确定，它们是
$\beta_{j,i}$-态射；这里

\VerdierSemanticAnchor{ch1-1-4-7-eq3}
\begin{equation}
  \beta_{j,i}(\sigma)=s_j(\sigma)s_i(\sigma)^{-1},
  \qquad i,j\in\{1,2,3\},\quad i<j.
  \tag{1.4.7.3}
\end{equation}

因此有关系式

\VerdierSemanticAnchor{ch1-1-4-7-eq4}
\begin{equation}
  \beta_{3,1}=\beta_{3,2}\beta_{2,1}.
  \tag{1.4.7.4}
\end{equation}

\VerdierSemanticAnchor{ch1-1-4-7-composition}
关系式 (1.4.7.2) 与 (1.4.7.4) 使我们能够定义 $\beta_{3,2}$-态射 $v$
与 $\beta_{2,1}$-态射 $u$ 的\emph{复合}；如此定义的复合与 $H$-函子态射的
复合相容。

\VerdierSemanticAnchor{ch1-1-4-8}
\noindent\textbf{1.4.8.}\quad
设 $\mathcal C'$ 为加性范畴，并设
\[
  (F,p)\colon(\mathcal C,G,T,c)\longrightarrow
  (\mathcal C',G,T',c')
\]
为一个由正规化余循环
$\alpha\in Z^2(G,\mathbb Z/2\mathbb Z)$ 扭曲的 $G$-函子。记

\VerdierSemanticAnchor{ch1-1-4-8-eq1}
\begin{equation}
  (F^{\circ},p^{\circ})\colon
  (\mathcal C^{\circ},G^{\circ},T^{\circ},c^{\circ})
  \longrightarrow
  (\mathcal C'^{\circ},G^{\circ},T'^{\circ},c'^{\circ})
  \tag{1.4.8.1}
\end{equation}
为函子 $(F,p)$ 的\emph{反扭曲分次函子}（沿用 (1.3.9) 的记号）。它定义如下。

\VerdierSemanticAnchor{ch1-1-4-8-opposite-data}
底层函子是 $F^{\circ}$，即 $F$ 的反函子。对每个
$g\in G^{\circ}$（$G$ 的反群），$\mathcal C^{\circ}$ 中的同构
\[
  p^{\circ}(g)\colon
  F^{\circ}T^{\circ}(g)=FT(g^{-1})
  \overset{\sim}{\longrightarrow}
  T'^{\circ}(g)F^{\circ}=T'(g^{-1})F
\]
等于 $p^{-1}(g^{-1})$。

\VerdierSemanticAnchor{ch1-1-4-8-opposite-cocycle}
函子 $(F^{\circ},p^{\circ})$ 是一个由余循环
$\alpha^{\circ}\in Z^2(G^{\circ},\mathbb Z/2\mathbb Z)$ 扭曲的
$G^{\circ}$-函子，其中 $\alpha^{\circ}$ 是 $\alpha$ 的反余循环：
\[
  \alpha^{\circ}(g,g')=\alpha^{-1}(g'^{-1},g^{-1}).
\]

\VerdierSemanticAnchor{ch1-1-4-9}
\noindent\textbf{1.4.9.}\quad
本段使用群 $\mathbb Z/2\mathbb Z$ 的唯一性质，是它典范地作用于加性范畴的
恒等函子。因此，可以用一个交换群 $K$ 代替 $\mathbb Z/2\mathbb Z$ 来发展
类似的讨论，只要 $K$ 以与 $G$ 的作用相容的方式作用于所考虑范畴的恒等函子。
